% C. Objetivos (Guía apdo. 3.C)
% Extensión orientativa: 2-3 páginas
\chapter{Objetivos}
\label{cap:objetivos}

\section{Objetivo general}

Diseñar e implementar un sistema modular de monitorización que emplee técnicas de
aprendizaje automático para el diagnóstico preventivo de sistemas electromecánicos,
analizando variables físicas --vibración, temperatura y sonido-- con el fin de identificar
desviaciones respecto al funcionamiento nominal del activo, y desplazando la lógica de
diagnóstico al extremo de la red.

\section{Objetivos específicos}

Los objetivos específicos se corresponden con las tareas comprometidas en la solicitud de
asignación de tema del presente TFM:

\begin{enumerate}
    \item \textbf{Definición del escenario de prueba.} Identificar un activo
        representativo y caracterizar su firma operativa.
    \item \textbf{Diseño de la arquitectura del sistema.} Definir el modelo de intercambio
        de datos entre los nodos sensores y la plataforma de monitorización, priorizando la
        eficiencia en el procesamiento local.
    \item \textbf{Selección de componentes de hardware y software.} Elegir la plataforma de
        procesado y los sensores adecuados para la captura de señales con la precisión
        requerida.
    \item \textbf{Desarrollo del motor de análisis inteligente.} Implementar algoritmos de
        aprendizaje automático para la extracción de características y la clasificación
        automática de los estados de salud del sistema monitorizado.
    \item \textbf{Validación experimental.} Realizar pruebas en un entorno controlado para
        calibrar el sistema y verificar la precisión del diagnóstico ante distintos tipos
        de fallo.
    \item \textbf{Pruebas en entorno real.} Desplegar una prueba de concepto sobre un
        activo operativo para validar la robustez de la arquitectura en condiciones de uso
        diario.
\end{enumerate}

\section{Alcance y limitaciones}

El trabajo se limita a dos activos de prueba: un compresor de referencia en estado nominal, sobre
el que se ajusta el modelo y se despliega el nodo, y un segundo compresor con una avería
preexistente que constituye el conjunto de evaluación. De ello se sigue la limitación de mayor
peso, desarrollada en el capítulo~\ref{cap:resultados}: al proceder la evidencia de detección de
una \textbf{máquina distinta} de aquella con la que se entrenó, el contraste no separa el efecto
del fallo del de la variabilidad entre ejemplares. Sobre el activo de referencia solo se
indujeron perturbaciones \textbf{operativas y ambientales} —reversibles y sin daño para la
máquina—, no fallos mecánicos. No se persigue tampoco la certificación metrológica de las
medidas: los sensores empleados son de bajo coste y las magnitudes se utilizan como indicadores
relativos de desviación, no como valores absolutos calibrados.

Procede acotar el alcance del motor de análisis en los términos en que efectivamente se ha
desarrollado. La propuesta de trabajo enuncia la clasificación automática de los estados de
salud; lo abordado es la \textbf{detección de una clase}, es decir el ajuste de un modelo del
estado nominal y la señalización de cuanto se desvíe de él, sin atribución de la desviación a un
estado concreto. La decisión responde a que un clasificador requiere ejemplos etiquetados de
cada estado, y los modos de fallo de un activo de esta naturaleza son escasos, variados y no
conocidos de antemano. La justificación detallada figura en el capítulo~\ref{cap:resultados}.

Tampoco se persigue el empleo de un entorno de ejecución de modelos reducidos para
microcontroladores. Esos entornos operan sobre redes neuronales, y el tamaño de muestra
alcanzable en una campaña de esta duración no permite ajustar una arquitectura de ese tipo sin
sobreajuste. La inferencia en el nodo se programa por tanto de forma directa, lo que no supone
una limitación del alcance sino una consecuencia de la familia de modelos aplicable.
